\section{Packets, finite horizons, and the induction invariant}\label[appendix]{sec:packets}

This appendix gives the geometry and bookkeeping for a finite number of
partitions.  The minimization, shielding, and regularization arguments
in the following appendices construct the regular cuts to which this
geometry applies.  The passage to countably many labels is proved in
\cref{countable:construction}.

The packet diagrams, finite stage constants, analytic invariant, and
radial initialization are stated in \Cref{sec:finite-cuts}. We prove here
their geometric properties and the compatibility of finite reconstructions.

\subsection{Product ends and regular splitting}

\begin{lemma}[Compatible splitting and finite truncations]\label{packet:split}
Suppose a packet diagram after the radial stage has the properties in
\cref{packet:category}.  On every packet let $g_E$ be a smooth compatible
function, with $0$ a regular value, including on every face.
Keep $L_E=g_E^{-1}(0)$ as the new hit space.  To each closed side
$\{\pm g_E\ge0\}$ attach $L_E\times[0,\infty)$ along its boundary.
Use smooth normal collars across the attachment, chosen as products
in all sufficiently far old tails.  Connected components of these
spaces, with the appended statuses, form a new packet diagram.

It has finite face and coface incidence, properly embedded closed tails,
and compact remainders.  Each packet admits a compact truncation which
is a smooth manifold with corners and is a strong deformation retract
of the full packet.  Truncations can be chosen simultaneously on every
finite downward diagram, with their faces equal to the corresponding
truncated packets.  If thresholds are in the product range of $g_E$,
the old retractions preserve its zero set and both closed sides.
These conclusions also hold at the first radial stage provided that,
before attachment, its wall and closed-side components are compact and
form a locally finite family in $X$.
\end{lemma}

\begin{proof}
First work in the finite downward diagram of one old packet.  Raise
all entry thresholds until every overlap is in a commuting quadrant,
and until all the functions in this diagram are products there.
For an elementary normal coordinate $t_i$ choose a larger horizon $h_i$.
On its tail use
$t_i\mapsto\min(t_i,h_i)$ and the homotopy
$t_i\mapsto(1-s)t_i+s\min(t_i,h_i)$; use the identity elsewhere.
The formulas agree continuously at $t_i=h_i$.  In an overlap the
other normal coordinates are unchanged and the formulas commute.
One may therefore perform these clamps in any order.  Their joint image
is compact: remove the far elementary tails, and then inductively
truncate each of the finitely many face tails.  This gives a finite
union of compact products.  Its boundary strata have precisely the
claimed product descriptions.  Each clamp fixes the cutting function,
since that function is a face pullback throughout the moving region.

The regular zero set in this compact manifold with corners is transverse
to every corner stratum.  Indeed on a far old face the derivative in
face directions is the derivative of the regular face function.
Consequently the compact zero set and the two compact closed sides are
manifolds with corners.  Each has finitely many components: every point
has a connected neighborhood in the relevant manifold with corners,
and compactness supplies a finite such cover.  Clamping shows that a
full zero or side component has a nonempty connected intersection with
the truncation.  Conversely a connection in the full set clamps to a
connection in the truncation.  Thus no additional components or
connections occur only at infinity.

On an old tail the whole operation is literally the same operation on
its face times the old normal intervals.  In particular the cylinder
attached to a face zero set is the same subset in every order of
passing to that face.  For the new sign-to-hit relation the tail is
exactly the newly attached cylinder.  These observations prove all
product and commutation rules, including the unique occurrence for
a prescribed coarser status: the outgoing product of a connected face
packet is connected and belongs to a unique coarser component.
A new packet has finitely many faces by the compact truncation just
constructed.  A face has at most one incident packet for each coarser
status, of which there are finitely many.  Applying these assertions
through the finite status depth proves finiteness in both directions.

For clarity, properness of tails is part of this argument, not an
inference from a metric completeness assertion.  Fresh closed cylinder
tails are proper by the attachment topology.  The continuous collapse
of a new space to its old parent sends a compact set to a compact set.
On an inherited tail, old properness then bounds its old parameters
in a compact subset of the old tail; the fresh cylinder length is
also bounded on the given compact set.  The remaining zero and side
conditions are closed.  Thus the inverse image of that compact set
in the inherited closed tail is compact.  Equivalently, it is contained
in the union of closed old pieces and bounded fresh cylinders over
compact zero pieces.  This also shows that the complement of all new
far tails is compact.  Clamp the fresh cylinder coordinate as well
to obtain the asserted new compact truncations and retractions.

At the radial stage there are initially no old tails.  The hypotheses
give compact wall and side pieces, with finite incidence on each
compact piece, and the fresh cylinders start the preceding induction.
All diagrams remain countable, because all manifolds involved are
second countable and the radial family is locally finite.

For later metric uses, normal collars can be obtained from any
compatible smooth parent metric.  In an old product tail, both the
regular defining function and the metric are pulled back from the
face, so the unit normal and its normal geodesics split off the old
normal coordinates.  On a compact truncation of the zero set choose
a positive tubular radius by finitely many normal-coordinate charts.
Induction through its finite face diagram extends this radius to a
positive uniform radius on that parent packet: any failure would
be in a product over a face where the radius was already chosen.
Thus these collars are products on all old ends; at the radial step
the same choice is made separately on each incident compact block.
\end{proof}

\subsection{Completeness, growth, and product limits}

\begin{lemma}[Metric geometry of a packet]\label{packet:geometry}
Equip a packet diagram after the radial stage with compatible smooth
base metrics and finitely many compatible smooth positive warp
functions.  Each connected packet $E$, separately, is complete and
has bounded geometry: positive injectivity radius and bounds on
curvature and every covariant derivative.  Every warp and all its
covariant derivatives are bounded on $E$, and every warp has a positive
lower bound.  For $\omega_E$ the product of the warps there are constants
$A_E<\infty$ and an integer $d_E\le\dim E$ such that
\begin{equation}\label{packet:growth}
 \int_{B_E(x,R)}\omega_E\,\dd\vol_E
       \le A_E(1+R)^{d_E}
       \qquad(x\in E,\ R\ge0).
\end{equation}
The constant is uniform over centers in this packet.  The same assertions
hold on $E\times\R^a$, with the product metric, pulled-back warps,
and growth exponent at most $d_E+a$, for every finite $a$.
\end{lemma}

\begin{proof}
Induct upward from packets with no faces, which are compact.
A full elementary product tail over a complete face has no finite-distance
exit through its face coordinate.  More precisely, a finite-length
curve in that product whose normal coordinate stays above the entry
level has Cauchy face and normal coordinates.  Completeness of the face
and properness of the closed tail keep its limit in the product.
It follows that, after fixing a closed-tail entry level $T_i$, the
function
\[
 \rho_i=\max(t_i-T_i,0)\quad\hbox{on that tail},\qquad
 \rho_i=0\quad\hbox{outside it}
\]
is continuous and globally $1$-Lipschitz.  The only exit is through
$t_i=T_i$, and its differential has norm at most one wherever defined.
A Cauchy sequence has all the finitely many $\rho_i$ bounded.
The compact-truncation assertion of \cref{packet:split} puts such a
sequence in a compact set.  Hence the packet is complete.

Recursively subdividing each coordinate into a bounded interval and
its far range gives a finite cover of $E$ by compact sets and sectors
\begin{equation}\label{packet:sector}
 K_B\times[T_1,\infty)\times\cdots\times[T_d,\infty),
\end{equation}
where $K_B$ is a compact subset of a further face and the metric and
warps are products.  The recursion terminates because a face decreases
base dimension.  On such a sector each displayed normal coordinate is
also the corresponding elementary coordinate on $E$, by commutation.
In an $R$-ball its range therefore has length at most $2R$, by the
Lipschitz functions $\rho_i$ (enlarge by a fixed bounded interval at
an entry if necessary).  The sector's weighted volume in that ball
is bounded by its compact base volume times a constant multiple of
$(1+R)^d$.  Summing the finitely many bounds proves
\cref{packet:growth}, for all centers, including centers outside the
sector.  Notice that $d\le\dim E$.

Bounds on smooth data follow from the same finite sector decomposition:
on compact pieces they are finite, and on a product they are precisely
face bounds.  For bounded geometry one can equivalently use uniform
charts.  Any failure of a uniform chart radius or of a fixed-order
chart bound has an escaping sequence; a subsequence lies in a far
product over a compact part of a further face, as detailed in
\cref{packet:escape}.  Fixed-radius neighborhoods there are product
neighborhoods with uniform smooth chart bounds and a positive
injectivity radius.  Compact remaining regions have the same
properties by a finite chart cover.  This proves bounded geometry
and all warp bounds.  The product assertions follow either by the
same proof or from
$B_{E\times\R^a}((x,z),R)\subset B_E(x,R)\times B_{\R^a}(z,R)$.
\end{proof}

\begin{lemma}[Escaping sequences and product windows]\label{packet:escape}
If a sequence escapes all compact sets of a fixed packet $E$, a
subsequence has coordinates $(x_\nu,t_{1,\nu},\ldots,t_{d,\nu})$
in one further-face product, where $d\ge1$, $x_\nu$ converges in a
compact subset of the face $B$, and every $t_{i,\nu}\to\infty$.
After translating these normal coordinates, the product charts contain
each fixed metric neighborhood of the sequence for all sufficiently
large $\nu$.  They exhaust $B\times\R^d$ locally.  All compatible
smooth data equal the corresponding product data on these windows.
\end{lemma}

\begin{proof}
Choose a sector from the finite cover in
\cref{packet:sector} containing a subsequence.  Pass to a subsequence
so each of its normal coordinates either converges finitely or tends
to infinity.  The compact base coordinate also converges.  Absorb the
bounded normal coordinates into the base of the intermediate face
specified by the remaining, unbounded coordinates.  Commutation
identifies this with the corresponding full product occurrence.
At least one unbounded coordinate remains, since otherwise the
sequence lies in a compact truncation.  A path of length at most $R$
changes each normal coordinate by at most $R$.  It cannot escape
through a face coordinate, by completeness of the face as in
\cref{packet:geometry}.  Thus a fixed-radius neighborhood lies in the
product once all its normal coordinates exceed the entry levels by
more than $R$.  Translating yields the stated exhaustion.  This
argument uses only the completeness part of \cref{packet:geometry},
which was established before its bounded-geometry conclusion.
\end{proof}

The initial metric and warps need not have bounded geometry on $X$.
At the radial split each bulk piece and each wall piece is compact
before cylinder attachment.  The smooth data have finite derivative
bounds there, the positive warps have positive minima, and the
cylinders freeze those data.  This initializes \cref{packet:geometry}.
For a locally finite initial warp list, a compact radial block and its
compact buffer meet only finitely many nonunit factors; the same
argument applies after omitting the unit factors.  Polynomial growth
is a conclusion for a single post-radial packet, never an assumption
on the original complete manifold.

\subsection{Compatible choices}

\begin{lemma}[Compatible choices and finite tolerance allocation]\label{packet:approximation}
On a packet diagram after the radial stage, choices of smooth cutoffs and approximations can
be made from deeper faces to cofaces, retaining exact products outside
finite thresholds.  More precisely, fix a finite derivative order and,
on each packet, positive error tolerances for any finite collection of
smooth data or approximations.  A datum on a face may additionally be
required to meet the tolerances of all its incident cofaces, including
through chains.  These requirements can be replaced in advance by
strictly positive face tolerances, and then met by compatible blending,
provided the individual local approximations and their required
convergence to the face data are available.
In particular there exist compatible smooth positive functions
arbitrarily small in any prescribed finite derivative norm on each
packet.  Also, if compatible closed sets $F_E$ and compatible open sets
$O_E$ satisfy $F_E\subset O_E$ on every packet, then
$\dist_E(F_E,E\setminus O_E)>0$ on each fixed packet whenever both
sets are nonempty.
\end{lemma}

\begin{proof}
There are only finitely many requests on a particular face, because
there are only finitely many coface chains.  Starting with the given
requests, propagate stricter tolerances down these chains, allowing
for the finitely many product, derivative, and cutoff constants.
Take finite minima; no minimum over an infinite adjacency component
is taken.  This finishes all tolerance choices before construction
begins.  A countable collection of packets causes no problem, since
each one has finitely many faces and cofaces and the depth is finite.

For the construction on one packet, enumerate its elementary tails.
Choose smooth cutoffs $f_i(t_i)$ rising from zero to one sufficiently
far out, with bounded derivatives through the chosen order, and
extend them by zero off their tails.  Write $u_0$ for an individual
approximation and $u_i$ for the already chosen face copy.  The weights
\begin{equation}\label{packet:blend}
 w_0=\prod_i(1-f_i),\qquad
 w_i=f_i\prod_{a<i}(1-f_a),\qquad
 w_0+\sum_iw_i=1
\end{equation}
produce $w_0u_0+\sum_iw_iu_i$.  Put the transitions beyond all relevant
intersection thresholds.  On an overlap the participating face copies
then agree by their own product rules.  In particular, wherever some
$f_i=1$, every nonzero term is the same face copy.  The result is an
exact copy there.  Leibniz' rule bounds the derivative error by a finite
multiple of the preassigned errors.  A perturbation on the remaining
compact region changes none of these end identities.
This proves the conditional approximation assertion.  It does not
assert convergence for a temporary nonproduct object; whenever such
an object is used below, its convergence is proved separately.

For positive small functions start with positive constants on the
compact deepest packets, take positive small individual choices on
compact parts, and use the positive weights above.  Decrease the
preassigned tolerances to cover all coface derivative requests.
This gives strict positivity at every point and any stipulated finite
norm bounds, including the $C^2$ bounds used for seam widths.
For the separation assertion, a sequence of distances tending to
zero either stays in a compact region, where closedness and inclusion
in the open set give a contradiction, or escapes.  In the latter case
\cref{packet:escape} supplies a common product window over a compact
part of a deeper face.  The two sequences project to the same limit
point, lying both in $F_B$ and in $B\setminus O_B$, a contradiction.
\end{proof}

\subsection{Finite reconstruction}

\begin{proposition}[Finite-horizon reconstruction]\label{packet:reconstruction}
Starting with $X$, perform finitely many compatible regular splits as
in \cref{packet:split}.  For the first split require locally finite
compact radial wall and side components.  At the later split for
color $i$, require the defining functions to be positive on the inverse
image of $A_i$ and negative on the inverse image of $B_i$, with those
inverse images contained in open regions of the stipulated signs.
After any $r\ge1$ of these splits, the diagram has a reconstruction
$X_r$, obtained by finite collar truncations and gluing at hit faces,
with the following properties:
\begin{enumerate}[label=(\roman*),leftmargin=*]
\item $X_r$ is a smooth manifold homeomorphic to $X$.  Its $r$ labelled
walls are closed smooth hypersurfaces with simultaneous transverse
intersection charts.  Each wall component is compact, the component
families are locally finite, and the closed sign/hit pieces are exactly
the compact cornered truncations of the analytic packets, by components.
\item The packet maps glue to a proper continuous map $q_r:X_r\to X$.
There is a homeomorphism $H_r:X_r\to X$ such that
\begin{equation}\label{packet:displacement}
 d_0\bigl(H_r(z),q_r(z)\bigr)\le r\qquad(z\in X_r).
\end{equation}
\item Horizons can be enlarged to exceed any already specified
compatible product thresholds on the full analytic packets.  A horizon
is constant for all occurrences of a normal coordinate belonging to
one connected global wall component.  Such enlargement has an
identification with the previous reconstruction which preserves
$q_r$, every old wall, and every old side status.
\item Once color $i$ has been processed, $q_r^{-1}(A_i)$ remains in its
open positive phase and $q_r^{-1}(B_i)$ remains in its open negative
phase at every subsequent stage.  In particular its wall misses
$q_r^{-1}(A_i\cup B_i)$.
\end{enumerate}
\end{proposition}

\begin{proof}
We prove the assertions in their order of dependence.  At stage zero
use $X_0=X$ and $q_0=H_0=\Id$.  At the radial split, the assumed compact
pieces give the first finite reconstruction by inserting cylinders
of finite length on the two sides of each wall component.  The
properness and homeomorphism argument below applies to this locally
finite compact family.  Its compact pieces also initialize the
analytic diagram by \cref{packet:split}.

Suppose the reconstruction through stage $r$ has been made.  Fix all
product thresholds for the next cutting functions and their collar
metrics in the \emph{full analytic coordinates} first.  Consider a
connected component $W$ of an existing global wall.  A packet with
the corresponding hit status has a nonempty connected truncation,
by the commuting retractions of \cref{packet:split}; it therefore lies
in a single such $W$.  The same assignment holds at all incidences
through further faces, since the corresponding truncations glue along
those faces.  Compactness of $W$ and local finiteness of the reconstructed
pieces imply that only finitely many such pieces meet $W$.  Thus there
are only finitely many threshold requests for its normal coordinate.
Choose one horizon exceeding all of them, with a positive buffer.
Do this for each existing wall component and coordinate.

This is a direct choice from previously fixed finite lists.  Increasing
a horizon cannot introduce a new analytic threshold: the functions
and the incidence diagram were fixed on the full infinite tails.
On a tail where topology and $q_r$ are already products, choose an
increasing smooth reparameterization from the enlarged finite normal
interval to the old one, equal to the identity before the old product
range and sending the new endpoint to the old endpoint.  Choose the
same reparameterization at all occurrences of $W$, on both sides of
its central hit wall.  It preserves those sides and the hit wall.
At corners the separately labelled normal reparameterizations commute.
Their product gives an identification
$\Phi_r:\widehat X_r\to X_r$ preserving every wall and sign status.
The glued position maps satisfy
\[
 \widehat q_r=q_r\circ\Phi_r,
 \qquad \widehat H_r=H_r\circ\Phi_r,
 \qquad
 d_0(\widehat H_r,\widehat q_r)\le r,
\]
because each position map is independent of the moved normal
variables.  Properness is preserved by this homeomorphism.  Thus
horizon enlargement introduces no error or change of position data.
For the rest of the step, relabel these enlarged representatives as
$X_r,q_r,H_r$.

The compatible cutting functions now glue to a single smooth function
on this enlarged $X_r$.  Near any old wall stratum the function is a
pullback from the hit piece, independent of all its signed normal
coordinates.  Its zero set $L$ is therefore regular and transverse
to every old wall stratum.  More explicitly, in a chart
$(b,s_1,\ldots,s_d)$ at an old stratum, the old walls are $s_j=0$
and the new function is $g_B(b)$.  The face hypothesis says
$\dd g_B\ne0$ at its zeros.  Hence the old normal covectors
$\dd s_1,\ldots,\dd s_d$ together with $\dd g_B$ are independent.
This also covers a stratum involving the radial wall.  At this point
$L$ has been formed without choosing a new horizon.  For a color split its forced signs give
\[
 q_r(L)\subset X\setminus(A_i\cup B_i).
\]
In the present one-ball-per-label setup,
$L\subset q_r^{-1}(B_{d_0}(x_i,2D))$.  The latter set has compact
closure because $q_r$ is proper and the original closed outer ball
is compact.  Since $L$ is closed in $X_r$, the whole color wall $L$
is compact.  In particular each component is compact, and there are
only finitely many components.  More generally, a closed regular
hypersurface with compact components has a locally finite component
family: a regular hypersurface chart meets only its own component,
and compact subsets of the hypersurface have finite such covers.
The first radial wall need not be compact as a whole; its stated
locally finite compact-component hypothesis gives exactly what is
needed for the same construction.

Only now assign a finite new horizon to each component of $L$.
Cut along $L$, attach a cylinder on each side, and identify the two
outer cylinder boundaries as the new central hit wall.  Near an old
wall this is the product of the operation on its hit face with the
old signed normal coordinate.  The same statement holds at multiple
old intersections.  Thus old walls extend smoothly through the new
cylinders, and all walls have transverse intersection charts.
The resulting pieces are exactly the truncations of the newly
appended statuses: the old clamps preserve $L$ and its signs and
extend across fresh cylinders keeping their new coordinate fixed;
the fresh-coordinate clamp commutes with all of them.  In particular
every full new packet retracts onto one nonempty connected piece
in the reconstruction.  This verifies the component assertion
needed to assign its coordinate at the next stage.

Let $c:X_{r+1}\to X_r$ collapse the inserted cylinders to $L$ and be
the identity on the bulk.  This map is proper.  A compact subset of
$X_r$ meets only finitely many new wall components; its inverse image
is contained in its compact bulk copy together with finitely many
finite cylinders over compact subsets of these compact components.
The glued map is $q_{r+1}=q_r\circ c$, and so is proper.
Every old compact wall meets finitely many new components, and its
transverse intersection with them is compact.  Its extension adds
only finitely many compact finite cylinders; the old wall therefore
remains compact.  Local finiteness of all pieces follows locally
from the finite collection of transverse wall charts, together with
the compact truncation descriptions.  This proves (i) and the
properness assertion without using any later wall.

We next choose collars that preserve all old strata.  The compatible
packet metrics glue on the enlarged reconstruction: near an old hit
stratum they have the form $h_B+\sum_j\dd s_j^2$.  On a neighborhood
of $L$, write $u$ for the glued cutting function and use the transverse
field
\[
 V=\frac{\grad u}{|\grad u|^2},\qquad V(u)=1.
\]
In each old product chart, $u=g_B(b)$ and this field has no old normal
component.  Thus it is tangent to every old wall meeting that chart
and is the pullback of the corresponding face field.  Equivalently,
one can patch such local fields with a compatible partition of unity:
$V(u)=1$ and tangency are linear conditions, so they survive patching.
The flow supplies product collars at old intersections.  It preserves
every old wall and each side, by uniqueness of the flow and continuity
from the identity.  This explains why a transverse ordinary collar
chosen without the product condition would not suffice here.

Choose pairwise disjoint locally finite strips around the compact
components of $L$ using this flow.  A positive strip width exists on
each component by compactness; it may be decreased to keep strips
disjoint.  For each component, continuity of $q_r$ and compactness
permit the width to be so small that $q_r$-images of any two points
on the same flow segment have $d_0$-distance less than one.  A monotone normal
reparameterization compresses the two finite inserted intervals and
the old strip into that old strip, fixing its outer boundary.  It
produces a homeomorphism $h:X_{r+1}\to X_r$, equal to the identity
off the strips, with
\[
 d_0\bigl(q_r(h(z)),q_r(c(z))\bigr)<1.
\]
Different strip widths and cylinder lengths are allowed on different
components.  Local finiteness makes these definitions a global
homeomorphism.  Set $H_{r+1}=H_r\circ h$.  The triangle inequality
and \cref{packet:displacement} give the new bound $r+1$.
The compression follows the stratum-preserving collar, so it also
preserves every old wall and its two side statuses.  The old signs
are pulled back under $c$, including across the cylinders, and the
newly imposed signs hold on the prescribed inverse images.  Pullbacks
of their open phases stay open.  This
proves (iv) and finishes the induction.
\end{proof}

\begin{remark}[Prescribed reconstruction errors]\label{packet:small-displacement}
The number $1$ in the last collar comparison can be replaced at stage
$r$ by any prescribed $\varepsilon_r>0$.  Each compact wall component
has its strip chosen separately, with the same error bound
$\varepsilon_r$ on every component.  Disjoint strips give a supremum
bound, rather than a sum over components.  The resulting error after
finitely many stages is at most $\sum_{j<r}\varepsilon_j$.
This freedom concerns the original position map and is independent
of the analytic lengths of the inserted finite collars.  The spatial
localization and limit of such reconstructions are treated in the
main reconstruction argument.
\end{remark}

\subsection{Initialization}

\begin{proof}[Proof of \cref{packet:initialization}]
The decreasing solution of
$2|y'|=.5y^2+\pi_0$, from $+\infty$ to $-\infty$, has finite
crossing length
\[
 \int_{-\infty}^{\infty}\frac{2\,\dd y}{.5y^2+\pi_0}
       =2\pi\sqrt{2/\pi_0}=T_0.
\]
Extend it by $+\infty$ before the crossing and by $-\infty$ after
it.  This is continuous as an extended-real-valued function, and
all its finite truncations are locally Lipschitz.  Composing the
transition, begun at distance $2$ from $A_i$, with the $1$-Lipschitz
function $d_h(A_i,\cdot)$ gives
\[
 2|\grad_h\mu_i|\le .5\mu_i^2+\pi_0
                    \le a_0\mu_i^2+\pi_0
\]
almost everywhere on the finite band.  Take
\[
 U_i^+=\{d_h(A_i,\cdot)<1\},\qquad
 U_i^-=\{d_h(A_i,\cdot)>T_0+3\}.
\]
Their closures lie strictly inside the respective pole interiors;
they contain $A_i$ and $B_i$.  The requested finite separations are
available from \cref{rounds:wide} after $N_*$, and hence $T_0$, are
fixed.  Their size does not change the original radius $D$.

We give the radial construction without a bounded-geometry
hypothesis.  Assume first that $X$ is noncompact, fix $o\in X$, and
put $d(x)=d_h(o,x)$.  Choose a countable locally finite cover by
relatively compact coordinate neighborhoods on which the metric
oscillation is small, and a subordinate smooth partition of unity
$\{\chi_j\}_{j\ge1}$ whose supports are compact in those neighborhoods.
Convolution in a slightly larger coordinate neighborhood gives
smooth $d_j$ with
\[
 |d_j-d|<\epsilon_j,\qquad |\dd d_j|_h\le\tfrac32
       \quad\hbox{near }\supp\chi_j.
\]
The derivative estimate follows from the local Lipschitz bound for
$d$ and the small metric oscillation; the convolution radius can
then be reduced to make its value error arbitrarily small.
Let $C_j=\sup|\dd\chi_j|_h<\infty$, and choose
$0<\epsilon_j<\min\{1/2,2^{-j-2}/(1+C_j)\}$.
The locally finite sum $\rho=2+\sum_j\chi_jd_j$ is smooth and
$|\rho-(2+d)|<1/2$.  Moreover, since $\sum_j\dd\chi_j=0$,
\[
 \dd\rho=\sum_j\chi_j\dd d_j
             +\sum_j(d_j-d)\dd\chi_j,
 \qquad |\dd\rho|_h\le\tfrac32+\sum_j2^{-j-2}<2.
\]
Completeness makes $d$ proper, so $\rho$ is proper and bounded below.
Thus a smooth proper radius with bounded gradient has been constructed
using only local smoothness and completeness.

Choose successive plateau intervals of the values of $\rho$, with
alternating signs and positive-width cores strictly inside each
plateau.  Separate successive plateaus by transition intervals of
width greater than $2T_0$.  On such an interval, a copy of the crossing
$y$, parametrized by $(\rho-a)/L$ with $L>2$ and with sign reversed
on alternate intervals, obeys the same profile inequality because
$|\dd\rho|_h/L\le1$.  Include an initial plateau extending below the
minimum of $\rho$.  Core closures have positive value margins to
both ends of their plateaus, except at the harmless lower end of the
initial plateau.  Sard's theorem supplies regular values in those
margins.  The sublevel domain at the initial plateau, and the domains
between two such regular levels at later plateaus, are compact smooth
barriers enclosing the corresponding core closures and lying in the
same pole interiors.  These barriers and the transition blocks form
locally finite families, since their value intervals tend to infinity.
Connectedness and properness imply that the image of $\rho$ contains
all sufficiently large values, so the separating subbands are nonempty.

A zero wall respecting the signs misses every separating core
subband.  On any connected wall component, the continuous function
$\rho$ therefore takes values in one bounded intervening block.
The component is closed in $X$, and properness of $\rho$ makes it
compact.  For a fixed sign, successive core subbands of the opposite
sign bound the $\rho$-values of every connected component of the
closed side; the central component is bounded by the first opposite
separator.  These closed-side components are compact as well.
In a compact block, regularity gives a finite cover by charts in
which the wall is a single sheet and each closed side is a half-space
or the full chart.  Only finitely many components can meet that block.
The blocks are locally finite, which proves the required local
finiteness of wall and side components.

If $X$ is compact, set the radial profile identically $+\infty$, with
$U_{\mathrm{rad}}^+=X$ and $U_{\mathrm{rad}}^-=\varnothing$.
The constant positive defining function realizes its empty wall;
the single nonempty closed side is $X$, which is compact.  All radial
profile and confinement assertions then hold directly.

Finally, the initial scalar requirement holds everywhere, since
\[
 \lambda_0+b_0\sum_{j\in J_0}\Pi_j
 =\frac14+1.12N_*\frac{1}{10N_*}
 =.362<.4\le R_X.
\]
The profile, core, and forced-preimage conditions have already been
verified.  This proves the invariant at stage zero.
\end{proof}

The geometric lemmas require compatible regular defining functions.
The minimizing boundaries used to obtain them may initially be
singular.  The remaining appendices first construct these minimizing
partitions, enlarge the remaining cores until the required guards
avoid the singular sets, and then produce the regular compatible functions to which
\cref{packet:split} applies.  The curvature restriction to the required
sets is respected precisely by \cref{packet:counts}.
